\section{评测策略综合对比}

综合前面五种模式，下表从多个维度进行对比，帮助读者根据具体场景选择合适的评测策略：

\begin{table}[H]
\centering
\caption{五大评测模式综合对比}
\small
\begin{tabular}{p{2.8cm}p{2cm}p{2cm}p{2cm}p{2.5cm}}
\toprule
\textbf{评测模式} & \textbf{执行成本} & \textbf{覆盖范围} & \textbf{调试难度} & \textbf{适用阶段} \\
\midrule
定制化测试逻辑 & 中等 & 视具体用例 & 低 & 全流程 \\
单步评测 & 低 & 单决策点 & 极低 & CI/CD 快速验证 \\
完整轮次评测 & 高 & 端到端 & 中等 & 发布前回归 \\
多轮对话评测 & 很高 & 用户交互 & 较高 & 关键场景覆盖 \\
环境隔离 & 基础设施 & 横切关注点 & --- & 全流程基础 \\
\bottomrule
\end{tabular}
\end{table}

\begin{importantbox}{选择评测模式的决策框架}
\begin{enumerate}
  \item \textbf{首先}确保环境隔离（模式五），这是所有其他模式的前提
  \item \textbf{主力}使用单步评测（模式二）覆盖关键决策点，约占用例的 50\%
  \item \textbf{补充}用完整轮次评测（模式三）做端到端验证
  \item \textbf{精选}几个关键场景用多轮评测（模式四）模拟真实交互
  \item \textbf{贯穿}使用定制化测试逻辑（模式一）为每个用例编写针对性断言
\end{enumerate}
\end{importantbox}


